% AUTO-GENERATED by scripts/build_wp_tables_v11.py -- do not edit by hand.
\begin{table}[htbp]
\centering\small
\caption{\textbf{Layer 1 point-in-time funnel and test statistics.} Firm counts use only filings submitted strictly before the as-of date.}
\label{tab:layer1-t}
\begin{tabular}{lrrr}
\toprule
& 2024-10-01 & 2025-04-01 & 2025-10-01 \\
\midrule
Filings on file & 5,932 & 6,866 & 9,508 \\
Firms on file & 3,575 & 3,575 & 3,577 \\
Firms with $\geq 2$ fiscal years & 2,357 & 3,288 & 3,574 \\
Firms screenable & 2,224 & 3,111 & 3,377 \\
Passing the quality gate & 87 & 153 & 176 \\
With a usable price series & 87 & 153 & 176 \\
Measurement start & 2024-09-30 & 2025-03-31 & 2025-09-30 \\
\midrule
\multicolumn{4}{l}{\textit{Newey--West $t$ on daily excess vs.\ TOPIX}} \\
Top 5, ROE rank & -1.10 & -1.92 & -1.28 \\
Top 20, ROE rank & -1.08 & -0.48 & -1.51 \\
Top 5, ROE${+}$margin rank & -1.46 & -0.28 & -0.98 \\
Top 20, ROE${+}$margin rank & -1.12 & -0.41 & -1.49 \\
\emph{Entire gate-passing pool} & -2.36 & -1.94 & -3.15 \\
\bottomrule
\end{tabular}
\end{table}

\begin{table}[htbp]
\centering\small
\caption{\textbf{Layer 3 (descriptive).} \emph{Gate-threshold stability of the mechanical selection.} Each variant perturbs exactly one threshold and re-runs the selection on the 2026 snapshot; there is no forward measurement here. Overlap counts how many of the 5 base holdings survive.}
\label{tab:robustness}
\begin{tabular}{lc}
\toprule
Variant (one threshold changed) & Overlap with base \\
\midrule
ROE threshold $\geq$12\% & 4 of 5 \\
ROE threshold $\geq$18\% & 5 of 5 \\
Operating margin $\geq$5\% & 4 of 5 \\
Equity ratio $\geq$40\% & 4 of 5 \\
Drop the revenue/profit-growth gate & 4 of 5 \\
Drop the loss-free gate & 4 of 5 \\
Rank on ROE alone & 4 of 5 \\
Sector cap raised to 3 & 5 of 5 \\
\midrule
\textbf{Minimum overlap across all 8 variants} & \textbf{4 of 5} \\
\bottomrule
\end{tabular}
\begin{minipage}{\textwidth}\vspace{2pt}\scriptsize
\textit{Notes.} Stability here is a property of the \emph{selection rule} on one snapshot, not evidence about returns. The mechanical selection is stable to single-threshold perturbation: at least 4 of 5 holdings survive every variant. This says the \emph{selection} is not knife-edge; it says nothing about whether the selection has forward-looking content, which Table~\ref{tab:layer1} addresses and answers in the negative.
\end{minipage}
\end{table}

\begin{table}[htbp]
\centering\small
\caption{\textbf{Attribution (hindsight-formed portfolio).} \emph{Cost sensitivity and the contest edition's daily convention.} Three-year annualised returns. The odd-lot spread of 15 bp is added to every one-way cost in the grid.}
\label{tab:cost-sens}
\begin{tabular}{lrrrrrrrr}
\toprule
& \multicolumn{5}{c}{Monthly rebalancing, one-way cost} & \multicolumn{2}{c}{Daily rebalancing} & Turnover \\
\cmidrule(lr){2-6}\cmidrule(lr){7-8}
Series & 0 bp & 10 bp & 25 bp & 50 bp & 100 bp & no cost & headline & (monthly) \\
\midrule
Portfolio (equal weight) & 53.3\% & 53.1\% & 52.8\% & 52.2\% & 51.2\% & 53.8\% & 50.9\% & 1.07$\times$ \\
Mechanical benchmark & 58.4\% & 58.2\% & 57.9\% & 57.3\% & 56.1\% & 58.7\% & 55.6\% & 1.15$\times$ \\
\bottomrule
\end{tabular}
\begin{minipage}{\textwidth}\vspace{2pt}\scriptsize
\textit{Notes.} The daily fixed-weight convention is the contest edition's and is not implementable; it appears here rather than in the main text. Even a 100 bp one-way assumption leaves the ranking of the two series unchanged, which is why Section~\ref{sec:costs} declines to treat cost robustness as a result.
\end{minipage}
\end{table}

\begin{table}[htbp]
\centering\small
\caption{\textbf{Measurement audit (no inference layer).} \emph{Evidence level against realised revenue growth --- the result runs against our own proposal.} The 14 hand-coded firms; annualised total revenue growth across the fiscal years on file.}
\label{tab:correspondence}
\begin{tabular}{lrrrr}
\toprule
Evidence level & $n$ & Median & Mean & Range \\
\midrule
Level 2 & 8 & +13.5\% & +14.5\% & -1.1\% to +28.3\% \\
Level 3 & 6 & +7.8\% & +7.5\% & -2.2\% to +21.3\% \\
\bottomrule
\end{tabular}
\begin{minipage}{\textwidth}\vspace{2pt}\scriptsize
\textit{Notes.} \textbf{Level-3 firms grew revenue more slowly than level-2 firms over this window.} We report this rather than omit it. It is not evidence that the ladder fails and we would resist a reader concluding either way: the cells hold 8 and 6 firms, the measure is \emph{total} revenue rather than theme-attributable revenue (which almost no firm discloses separately), the levels were not coded as-of, and nothing controls for sector or size. The table has no power to order the levels --- which is precisely the point made in Section~\ref{sec:ladder-validation}.
\end{minipage}
\end{table}

